%% ==========================================================================
%%  9  Conclusion and Future Work
%%
%%  Mirrors R3's Section 5.  The two "don't go here" results are the reason
%%  this section is worth reading before starting work.
%% ==========================================================================
\section{Conclusion and Future Work}
\label{sec:conclusion}

We have not settled Conjecture~\ref{conj:main}, and the honest summary of this
paper is that it makes the conjecture harder to prove than it looked. Two
strategies that a reader would reasonably try first are ruled out:
item-by-item insertion, which is the strategy that works
in~\cite{BKNS22} (Theorem~\ref{thm:obstruction}), and restriction to
utilitarian-optimal allocations (Proposition~\ref{prop:msw-false}). What remains
is a small set of positive results --- binary additive costs
(Theorem~\ref{thm:binadd}), identical costs (Theorem~\ref{thm:identical}), two
agents --- and three structural tools that between them say precisely how the
chore problem differs from its goods mirror.

The single most useful thing to keep in mind is the constraint the two negative
results jointly impose: \emph{a correct algorithm must be free to move
already-allocated chores between bundles, and free to give up total cost while
doing so}. Every construction in this paper that succeeds does so by exhibiting
a per-agent potential $\phi$ with $\edgew{\alloc}(i,j) \le \phi_i - \phi_j$,
which telescopes along paths. Finding such a potential for general negative
dichotomous costs, or proving none exists, is the problem.

\paragraph{Directions.}
\begin{enumerate}
  \item \textbf{Adversarial instance generation} (\Cref{sec:samplerbias}). The
    randomised evidence is drawn from a distribution that essentially never
    produces the structured extremal functions from which both negative results
    are built. This is cheap, and it is the only inexpensive way to find out
    whether Conjecture~\ref{conj:main} is true before investing in a proof.
  \item \textbf{Exchange paths rather than single transfers}
    (\Cref{sec:tiebreaks}). Single-chore local search has bad local minima. The
    move set should match the exchange structure of dichotomous costs, as
    Yankee Swap and matroid union do in the matroid-rank
    literature~\cite{BEF21}; and by Remark~\ref{rem:msw-deadend} the search must
    range over all allocations, not the utilitarian-optimal ones.
  \item \textbf{Build the two tiers first} (Proposition~\ref{prop:twotier}).
    Choose the partition into paid and unpaid agents and only then allocate.
    This is the one natural approach that Theorem~\ref{thm:obstruction} does not
    obviously kill, since it is not incremental in the chores. In the proof of
    Theorem~\ref{thm:binadd} the paid set is exactly the agents holding a
    $\lceil q/n \rceil$-share of the universally disliked chores; identifying
    what plays that role in general is the first question.
  \item \textbf{Climb the milestone ladder.} Binary additive is done
    (Theorem~\ref{thm:binadd}). Binary submodular is the next rung: for
    matroid-rank costs, matroid union should supply the cardinality balancing
    that Theorem~\ref{thm:sizeshift} identifies as the missing ingredient.
    General dichotomous is the summit.
  \item \textbf{The mixed model.} Allow each item to be a good or a chore for
    each agent, with all marginals in $\set{-1,0,1}$: the dichotomous
    restriction of the doubly monotone setting of~\cite{KMSTY24}. This is the
    natural common generalisation of~\cite{BKNS22} and of the present paper.
\end{enumerate}

\paragraph{Before investing further.}
The best published bound covering negative dichotomous valuations is
that of~\cite{KMSTY24}: $n-1$ per agent and $n(n-1)/2$ in total, from any EF1
allocation. Conjecture~\ref{conj:main} would improve the total by a factor of
$n$ on this subclass --- the same factor~\cite{BKNS22} gains
over~\cite{BDNSV20} on the goods side.
\todo{confirm that nobody has published ``dichotomous chores with subsidy''
already. This is the one check that could make the whole programme redundant and
it has not been done.}
